\chapter{Thiết kế và hiện thực hệ thống}
\label{chap:implementation}
\section{Yêu cầu hệ thống}
Hệ thống ART gồm Decoder tùy biến từ Qwen2, bộ truy hồi tài liệu, bộ phân đoạn câu trực tuyến và bộ điều phối huấn luyện. Bộ truy hồi sử dụng E5 để mã hóa truy vấn, kết hợp FAISS để tìm kiếm các đoạn văn liên quan. Quá trình mã hóa áp dụng tiền tố truy vấn E5, mean pooling theo attention mask và chuẩn hóa L2; các vector float32 được sử dụng để tìm kiếm trên CPU. Mỗi kết quả truy hồi được ánh xạ đến đoạn văn tương ứng trong corpus và cung cấp cho Decoder dưới dạng quan sát.

\section{Kiến trúc hệ thống}
\subsection{Chuẩn bị chuỗi và tập tham số}
Câu hỏi được định dạng theo mẫu hội thoại, trong khi chuỗi đích do teacher model sinh ra được nối với token EOS và mã hóa riêng với câu hỏi. Các token thuộc câu hỏi và nội dung tài liệu truy hồi được gán nhãn -100 để loại khỏi quá trình tính LM Loss. Bộ nạp dữ liệu kiểm tra kiểu chuỗi và loại bỏ các mẫu có chuỗi đích rỗng.

Checkpoint ART được sử dụng trong đánh giá được huấn luyện trong 0.75 epoch. Quá trình fine-tuning chỉ cập nhật 15 lớp Decoder cuối trong tổng số 36 lớp cùng với Decision Head; các lớp còn lại được đóng băng nhưng vẫn tham gia lan truyền thuận. Kích thước micro-batch là 1, với gradient accumulation qua 32 bước.

\section{Quy trình Collect--Replay và suy luận}
\Cref{fig:workflow-overview} trình bày tổng quan ba giai đoạn hoạt động của ART.

\par\medskip\noindent
\begin{minipage}{\linewidth}
\centering
\captionsetup{hypcap=false}
\includegraphics[width=0.75\textwidth,keepaspectratio]{graphics/art-retrieval-workflow-approved.png}
\captionof{figure}{Sơ đồ tổng quan quy trình hoạt động của ART.}
\label{fig:workflow-overview}
\end{minipage}\par

\subsection{Collect--Replay khi training}
Collect–Replay tách bước chọn quỹ đạo truy hồi rời rạc khỏi bước tối ưu có gradient.

\textbf{Collect.} ART nạp câu hỏi, rồi xử lý chuỗi đích theo segment. Tại các điểm ứng viên, mô hình so sánh điểm mới \(d_t\) với ngưỡng. Ở các biên câu đủ điều kiện, exploration cũng có thể kích hoạt truy hồi. Collect chạy trong \texttt{no\_grad}, lưu tiền tố, attention masks, điểm và ngữ cảnh theo quỹ đạo. Nội dung quan sát lấy từ bộ truy hồi được chèn vào tiền tố, căn chỉnh với nhãn bỏ qua.

\textbf{Replay.} Mô hình chạy lại các tiền tố đã lưu mà không dùng KV Cache, tính lại điểm quyết định trong đồ thị gradient và ghép các phần logits mới theo thứ tự chuỗi. Quá trình này giữ nguyên tài liệu đã thu thập, tính \cref{eq:loss} và backpropagate. Quỹ đạo có thể kết thúc bằng một bản ghi chỉ chứa quyết định; nhánh này có thể cộng điểm cuối lần nữa và không đồng nghĩa truy hồi.

\subsection{Inference tự hồi quy}
Ở prefill, ART khởi tạo ngữ cảnh và KV Cache. Trong các bước sinh sau, mô hình kiểm tra biên câu và chỉ tính điểm ứng viên tại các vị trí đủ điều kiện. Khi \(d_t\geq\tau\), hệ thống lấy câu hoàn chỉnh gần nhất hoặc câu hỏi gốc làm truy vấn và kiểm tra dấu hiệu bị chặn. Nếu chấp nhận, ART truy hồi, chèn quan sát, cập nhật document mask và đồng bộ cache nhánh thứ hai từ nhánh có quan sát. Điểm vừa chấp nhận trở thành hệ số trộn cho quan sát mới. Nếu không truy hồi, hệ thống giữ ngữ cảnh đang hoạt động và tiếp tục sinh. Inference không dùng exploration.

\Cref{fig:workflow} đối chiếu các nhánh quyết định ở Collect và inference, cùng đường truyền gradient trong Replay.
\par\medskip\noindent
\begin{minipage}{\linewidth}
\centering
\captionsetup{hypcap=false}
\includegraphics[width=\linewidth]{graphics/workflow-figure.pdf}
\captionof{figure}{Sơ đồ kỹ thuật quy trình Collect--Replay và suy luận của ART.}
\label{fig:workflow}
\end{minipage}\par
